\documentclass[a4paper,11pt]{article}
\usepackage[utf8]{inputenc}
\usepackage[T1]{fontenc}
\usepackage[french]{babel}
\usepackage[left=2cm,right=2cm,top=2cm,bottom=2cm]{geometry}
\usepackage{xcolor}
\usepackage{hyperref}
\usepackage{fontawesome5}
\usepackage{parskip}

% --- Couleurs Thales ---
\definecolor{primary}{RGB}{0, 56, 101}

\begin{document}

% --- EN-TÊTE ---
\begin{center}
    {\Large \textbf{Loïc Aron MBASSI EWOLO}}\\[4pt]
    \small
    \faMapMarker\ Cergy, France (Mobile France entière) \quad
    \faPhone\ (+33) 7 59 52 33 98 \quad
    \faEnvelope\ \href{mailto:loicaron.mbassiewolo@ensea.fr}{loicaron.mbassiewolo@ensea.fr}
\end{center}
\vspace{0.5cm}

% --- DESTINATAIRE ---
\hfill
\begin{minipage}[t]{0.45\textwidth}
    \textbf{Thales}\\
    Site de Laval\\
    Équipe FPGA -- Activité IFF\\
    Pays de la Loire, France
\end{minipage}

\vspace{0.8cm}

\textbf{Objet :} Candidature Stage Ingénieur Développement FPGA (MPSoC/RFSoC) -- Avril 2026\\
\textit{Réf : Stage Banc de Tests FPGA Zynq Ultrascale+}

\vspace{0.5cm}

Madame, Monsieur,

Votre offre de stage sur \textbf{RFSoC AMD/Xilinx Zynq Ultrascale+} pour l'élaboration d'un banc de tests IFF a immédiatement retenu mon attention. Actuellement en double diplôme ingénieur à l'\textbf{ENSEA} (Cergy) et à l'\textbf{École Polytechnique de Yaoundé}, je suis disponible dès \textbf{avril 2026} pour 4 à 6 mois.

Ma formation à l'ENSEA m'a permis d'acquérir les compétences clés pour votre projet. En \textbf{conception FPGA (VHDL)}, j'ai développé des descriptions matérielles, réalisé des simulations comportementales et effectué la synthèse via \textbf{Vivado}. En \textbf{programmation C bas niveau}, mon projet TAXI (simulation temps réel multiprocessus) m'a confronté à la gestion mémoire partagée, aux IPC et à la synchronisation par signaux UNIX -- des problématiques proches de l'intégration hardware/software sur MPSoC.

Je participe actuellement au \textbf{Challenge CoHoMa} (robotique autonome militaire, partenariat ENSEA/Armée française). Ce projet m'a permis de travailler sur des \textbf{plateformes embarquées Nvidia Jetson} avec des contraintes de performance et consommation similaires aux MPSoC, d'intégrer des capteurs complexes (Lidar VLP16, caméras profondeur) et de collaborer avec des équipes pluridisciplinaires pour assurer la cohérence fonctionnelle du système.

Mon expérience hardware est également concrète : lors du projet \textbf{Harmony Gloves} (gant connecté pour langue des signes), j'ai participé au prototypage complet incluant la \textbf{soudure} de composants, la programmation \textbf{STM32} et l'implémentation des protocoles \textbf{I2C/SPI}. Mon projet d'option sur le \textit{Fault Detection \& Fault-Tolerant Control} pour drone m'a familiarisé avec la modélisation de systèmes critiques sous \textbf{MATLAB/Simulink}.

Intégrer l'équipe FPGA du site de Laval représente une opportunité exceptionnelle : travailler sur des composants de dernière génération dans un contexte applicatif concret (radars IFF), contribuer à l'ensemble du cycle de développement jusqu'à la \textbf{validation sur banc de tests}, et apprendre aux côtés d'un architecte expérimenté et de 5 développeurs FPGA.

Rigoureux, autonome et doté d'un fort esprit d'équipe, je serais honoré de mettre mes compétences au service de Thales.

Je vous prie d'agréer, Madame, Monsieur, l'expression de mes salutations distinguées.

\vspace{0.8cm}

\textbf{Loïc Aron MBASSI EWOLO}\\
\textit{Élève-Ingénieur ENSEA / Polytechnique Yaoundé}

\end{document}
